% ---- pg-029 （书页 17） ----

\begin{mvfigure}
\mvinc{pg-029}{0.29975,0.13067,0.73326,0.73322}
\phantomsection\label{fig:1.11}\mvcap{图 1.11\quad 改进的散布图对(a)选择的点的组和(b)删去标本 25 在内的诸点,并重画散布图}
\end{mvfigure}

的透视图.这些视图是通过连续地旋转及转换三维坐标轴得到的.旋转坐标轴我们可得到对数据的三维外观的更好的了解.\hyperref[fig:1.12]{图 1.12}(d)给出硬度数据在 $x_3, x_1, x_2$ 空间的图片.注意到\hyperref[fig:1.12]{图 1.12}(a)和(d)清楚地证实了样本 9 和 16 是离群值.样本 9 的数值在所有的三维坐标中都非常大.把\hyperref[fig:1.12]{图 1.12}(a)的坐标轴逆时针旋转得到\hyperref[fig:1.12]{图 1.12}(b),两个异常的观测值在这个视图中被屏蔽了.进一步旋转坐标轴 $x_2, x_3$ 得到\hyperref[fig:1.12]{图 1.12}(c);此时离群值之一标本 16 隐藏起来.

通过这些缓慢旋转的数据进行目测有时可获得额外的认识.统计学家们对这种动态的方式刚刚才开始了解和利用.
